% from main.tex
%%%%%%%%%%  scan idx 115  |  folio 102  %%%%%%%%%%
monomorphisme induisant l'identité en degré \(\geq n\).

\textbf{Définition 2.8.} Soient \(n\in\mathbb N\) et \(F\in\Ob\mathcal C_S\). Nous dirons que \(F\) est de \(n\)-\(\mathcal C_0\)-présentation finie si l'ensemble des \(X\) au-dessus de \(S\) tels qu'il existe une suite exacte
\[
  E^{-n}\longrightarrow\cdots\longrightarrow E^0\longrightarrow F_X\longrightarrow 0,
\]
avec \(E^i\in\Ob\mathcal C_{0,X}\) pour \(-n\leq i\leq 0\), est un raffinement de \(S\). Nous dirons que \(F\) est d'\(\infty\)-\(\mathcal C_0\)-présentation finie si \(F\) est de \(n\)-\(\mathcal C_0\)-présentation finie pour tout \(n\in\mathbb N\).

Bien entendu, nous dirons parfois \(n\)-présentation finie tout court quand il n'y aura pas d'ambiguïté possible sur \(\mathcal C_0\). Nous dirons « de présentation finie » pour « de 1-présentation finie », et « de type fini » pour « de 0-présentation finie », cette dernière notion coïncidant avec celle déjà introduite en (1.2).

\textbf{Proposition 2.9.} Soient \(n\in\mathbb N\) et \(F\in\Ob\mathcal C_S\). Conditions équivalentes :

(i) \(F\) est de \(n\)-présentation finie (resp. d'\(\infty\)-présentation finie);

(ii) \(F\) est \((-n)\)-pseudo-cohérent (resp. pseudo-cohérent).

Dans le cas où \(S\) est un site ponctuel, \(F\) est pseudo-cohérent si et seulement s'il existe une suite exacte
\[
  \cdots\longrightarrow E^i\longrightarrow\cdots\longrightarrow E^0\longrightarrow F\longrightarrow 0,
\]
avec \(E^i\in\Ob\mathcal C_0\) pour tout \(i\).

\emph{Preuve.} Pour l'équivalence \((\mathrm{i})\Leftrightarrow(\mathrm{ii})\), il suffit évidemment de traiter le cas non respé. L'implication \((\mathrm{i})\Rightarrow(\mathrm{ii})\) est immédiate sur les définitions, et l'%
